\appendix
\chapter{The Framework and Its Empirical Tests}

The main text uses mechanisms to organize the gap between proposals and dependable change. This appendix specifies how to turn that organization into a measurement program.

\section{Intellectual Lineage and Proposed Contribution}

The distinction between bits and atoms is familiar from discussions of digitization, including Negroponte's \textit{Being Digital} \cite{negroponte1995being}. This book adds life as an analytical category because adaptive, self-maintaining systems can change the pathways through which an intervention acts.

Ogburn's cultural-lag account directs attention to differing rates of technical and institutional change \cite{ogburn1922social}. Productivity debates associated with Solow direct attention to the gap between available computation and measured economic outcomes \cite{solow1987watch}. Baumol and Bowen explain how productivity differences can change relative service costs \cite{baumol1966performing,baumol1967macro}.

These traditions contain richer arguments than any short genealogy can preserve. They identify related problems without establishing one common underlying law.

The proposed contribution is a shared diagnostic language: resource barriers $R$, temporal rigidity $T$, and topological coupling $G$. The framework follows a capability through design, execution, verification, operation, and responsibility. It asks whether identifying those constraints improves explanation and decisions across domains.

The framework has three levels. A classification organizes attention. A fitted quantitative model relates operational variables to outcomes. A normative argument concerns which results, rights, and responsibilities deserve protection. Each requires its own evidence.

A useful test is comparative. Does the framework identify a bottleneck earlier, improve prediction, reveal an omitted cost, or distinguish interventions that a simpler account treats alike? The value of the language lies in those gains.
